\section{Endpoint curvature of the actual theta kernel}
\label{sec:curvatureendpoints}
The condition $V'''(u)>0$ for $u>0$,
with $V=-\log\Phi$, remains globally unresolved.
We isolate both endpoint regimes by exact analytic arguments.

\begin{proposition}[Modular-point cancellation]
Let
\[
\Phi_1(u)=\pi(2\pi e^{4u}-3)e^{5u-\pi e^{4u}},
\qquad V_1=-\log\Phi_1.
\]
For $y=\pi e^{4u}$,
\begin{align}
V_1'''(u)
&=\frac{64y}{(2y-3)^3}P(y),\\
P(y)&=(2y-3)^3-6(2y+3).
\end{align}
In particular $V_1'''(0)<0$, whereas
$V_1'''(u)>0$ whenever $y\ge7/2$.
The polynomial has exactly one real root
on $[3,\infty)$, lying between $13/4$
and $7/2$.
\end{proposition}
\begin{proof}
Writing $\log\Phi_1(u)=
\log\pi+\log(2y-3)+5u-y$
and applying $(4y\,d/dy)^3$ yields
the displayed expression. The polynomial is
\[
P(y)=8y^3-36y^2+42y-45.
\]
Its derivative $P'(y)=24(y-3/2)^2-12$
is positive for $y\ge3$. The bounds
$3<\pi<22/7<13/4$ and $P(13/4)=-113/8<0$
give $P(\pi)<0$, while $P(7/2)=4>0$.
For $y=t+7/2$ with $t\ge0$,
\[
P(y)=8t^3+48t^2+84t+4>0.
\]
Monotonicity and continuity give uniqueness.
\end{proof}

\begin{proposition}[Full-kernel endpoint positivity]
Write $s(t)=\Phi(\sqrt t)=a_0+a_1t+a_2t^2+O(t^3)$
with $a_0>0$. If $\mathcal T s(0)
=a_1^2-2a_0a_2>0$, then
\begin{equation}
V'''(u)=\frac{12\,\mathcal T s(0)}{s(0)^2}
\,u+O(u^3)\quad (u\to0^+).
\end{equation}
Hence $V'''(u)>0$ for all sufficiently
small $u>0$ for the genuine theta kernel.
Moreover $V'''(u)>0$ for all sufficiently
large $u$, without assuming a global
curvature hypothesis.
\end{proposition}
\begin{proof}
The theta functional equation makes $\Phi$
even near the modular fixed point.
Expanding
\[
-\log s(u^2)=-\log a_0-\frac{a_1}{a_0}u^2
-\left(\frac{a_2}{a_0}
-\frac{a_1^2}{2a_0^2}\right)u^4+O(u^6)
\]
and differentiating three times yields
the first assertion. Strict positivity of
$\mathcal T s(0)$ is available from the
earlier second-level kernel theorem
and independently from the exact interval
certificate in this work.

For large $u$ put $y=\pi e^{4u}$. The ratio
of the full sum to its first summand is
$1+R(y)$, where
\[
R(y)=\sum_{m\ge2}
m^2\frac{2m^2y-3}{2y-3}e^{-(m^2-1)y}.
\]
Termwise derivatives through order three
are uniformly dominated by a polynomial
in $y$ times $e^{-3y}$ for $y\ge y_0>3$;
this follows by a convergent geometric
majorant after bounding all factors by
polynomials in $m$ and $y$.
It follows that
$(4y\partial_y)^3\log(1+R(y))
=O(y^N e^{-3y})$ for some fixed $N$.
Since the first-summand curvature above
satisfies $V_1'''(u)=64y+O(1)$, we obtain
\[
V'''(u)=64\pi e^{4u}+O(1)
\qquad(u\to\infty),
\]
hence eventual strict positivity.
\end{proof}

The global problem is now concentrated on a
\emph{compact intermediate interval}. The
modular neighborhood is controlled through
the positive second-level defect; the distant
tail is controlled by first-summand dominance.
An explicit rational subdivision certificate
for the intermediate region, followed by an
analytic/Lean transfer, is still missing.
This endpoint theorem does not imply the
complete Laguerre hierarchy on the real line,
nor the Riemann hypothesis.
